\section{Einführung in die Analysis I}

Die Analysis ist ein grundlegendes Teilgebiet
der Mathematik. Sie beschäftigt sich insbesondere
mit der Untersuchung von Funktionen, ihren
Eigenschaften und ihrem Verhalten.

Dieses Skript behandelt die wesentlichen
Grundlagen der Analysis I und führt schrittweise
in die wichtigsten mathematischen Methoden ein.

Ausgehend von Mengen, Abbildungen und der
vollständigen Induktion werden zunächst die
komplexen Zahlen sowie Folgen und deren
Konvergenz untersucht.

Darauf aufbauend folgen Grenzwerte und
Stetigkeit von Funktionen, elementare Funktionen,
Differentiation und höhere Ableitungen.

Weitere Schwerpunkte bilden die Integration,
die Untersuchung unendlicher Reihen und
Potenzreihen sowie die Taylorentwicklung.

Ergänzend wird die Fourieranalyse als
weiterführende Methode zur Darstellung
periodischer Funktionen behandelt.

Im Mittelpunkt stehen nachvollziehbare
Herleitungen, anschauliche Beispiele und
Übungsaufgaben. Den Abschluss bildet ein
Kapitel mit Aufgaben zur Wiederholung
und Vertiefung der behandelten Inhalte.